\section{Exact candidate-set computation and conditional direct formulations}
\label{sec:direct}

LexUR is primarily a candidate-set selection rule, and that is the setting of our
experiments. We first present the exact finite algorithm. We then develop, more tentatively, how the recommendation can
in principle be computed \emph{directly} on a continuous feasible set, replacing
the usual two-stage pipeline by a
single robust search. We present the continuous case
as a formulation with assumptions whose large-scale empirical validation we leave
to future work.

\subsection{Exact algorithm for finite candidate sets}

The finite-set procedure is summarised in Figure~\ref{alg:finite}.
\begin{figure}[h]
\centering
\fbox{
\begin{minipage}{0.9\textwidth}
\textbf{Procedure (Figure~\ref{alg:finite}): Exact finite-set LexUR ($\tau=0$) and tolerance reporting}\vspace{2mm}\\
\textbf{Input:} Candidate set $A=\{x_1,\dots,x_N\}$; objectives $F(x)$; labelled monotone probe family $\Q=\{q_1,\dots,q_K\}$; degeneracy threshold $\epsilon_q$.\\
\textbf{Output:} Exact minimiser set and labelled certificates; optionally a numerical $\tau$-class.\vspace{2mm}\\
\textbf{1. Dominance filter:} Remove any Pareto-dominated $x \in A$.\\
\textbf{2. Normalisation:} For each criterion $i$:
\begin{itemize}
    \item Compute active range $\ideal_i = \min_A f_i(x)$ and $\nadir_i = \max_A f_i(x)$.
    \item If $\nadir_i-\ideal_i>\epsilon_q$, compute $r_i(x)=(f_i(x)-\ideal_i)/(\nadir_i-\ideal_i)$; otherwise flag the criterion and require external bounds or exclude probes that use it.
\end{itemize}
\textbf{3. Anchor bounds:} For each probe $q \in \Q$:
\begin{itemize}
    \item Evaluate $V_{iq} = q(r(x_i))$ for all active $x_i \in A$.
    \item Find $q^\star = \min_i V_{iq}$ and $q^{\mathrm w} = \max_i V_{iq}$; flag and exclude a probe when $q^{\mathrm w}-q^\star\le\epsilon_q$.
\end{itemize}
\textbf{4. Regret evaluation:} For each $x \in A$ and each valid probe $q \in \Q$:
\begin{itemize}
    \item Compute $\reg_q(x) = (V_{iq} - q^\star) / (q^{\mathrm w} - q^\star)$.
\end{itemize}
\textbf{5. Certificate sorting:} For each $x$, sort the labelled pairs $(\reg_q(x),q)$ descending to obtain $\reg^{\downarrow}(x)$ without discarding probe names.\\
\textbf{6. Exact selection:} With $\tau=0$, return every row attaining the lexicographic minimum.\\
\textbf{7. Numerical reporting (optional):} Around the exact minimiser $x^\star$, report $W_\tau$ or the stability set $C_\tau(x^\star)$ defined in \S\ref{sec:core}; do not treat this tolerance relation as a preorder or call it an equivalence class.
\end{minipage}
}
\caption{Finite-set LexUR procedure.}
\label{alg:finite}
\end{figure}

The worst-case complexity is straightforward. With naive probe evaluation, winner selection costs $O(NKm + NK\log K + NK)$. Producing a complete lexicographic ranking of all candidates costs $O(NKm + NK\log K + NK\log N)$. Here $O(NKm)$ is the cost of probe evaluation and $O(NK\log K)$ is the cost of sorting the $K$-dimensional disappointment vector for each of the $N$ candidates. The number of probes $K$ is controlled by the construction of \S\ref{sec:probes}.

\subsection{Cost of the anchors}
Both the finite and continuous formulations need, for each probe, its attainable
best $q^\star=\min_x q(r(x))$ and worst $q^{\mathrm w}=\max_x q(r(x))$. On a
candidate set these are $O(N)$ scans; on a continuous $\X$ each is a scalar global
optimisation, so the family contributes $2K$ such subproblems. Linear probes on
a polyhedral feasible set yield linear anchor problems. For a general convex
probe, minimising the probe may remain tractable, but maximising it over a convex
set is not generally a convex optimisation problem. Non-convex feasible sets or
probes may require global optimisation. Anchor computation can therefore be a
substantial bottleneck and is not assumed cheaper than front approximation in
all continuous cases.

\subsection{Front-free formulation for structured cases}
By formulating LexUR directly on $\X$, a front-free computation is theoretically available for structured cases (e.g., linear multiobjective problems), as we demonstrate empirically in \S\ref{sec:exp}. Lexicographic minimisation of the sorted regret vector over a continuous set $\X$ is exactly equivalent to sequentially minimising the cumulative sums of the $p$ largest disappointments, $S_p(\reg(x))=\sum_{k=1}^p \reg_{[k]}(x)$. For $p = 1,\dots,K$, this admits an exact ordered-value linearisation \citep{ogryczak2006}:
\begin{equation}
S_p(\reg(x)) = \min_{\alpha \in \R, u \in \R^K} \left[ p\alpha + \sum_{q \in \Q} u_q \right] \quad \text{s.t.} \quad u_q \ge \reg_q(x) - \alpha, \quad u_q \ge 0 \quad \forall q \in \Q.
\end{equation}

Lexicographic minimisation is implemented sequentially:
\begin{align}
\gamma_1 &= \min_{x \in \X} S_1(\reg(x)), \\
\gamma_p &= \min_{x \in \X} S_p(\reg(x)) \quad \text{s.t.} \quad S_j(\reg(x)) \le \gamma_j \quad \forall j < p.
\end{align}

If $\X$ is polyhedral and the probes are linear, each stage $p$ can be solved as an LP (or MILP, if $\X$ contains integer variables) over the decision variables $x$, the scalar $\alpha$, and the $K$ auxiliary continuous variables $u_q$. Max-linear probes (e.g., $q(r) = \max_i w_i r_i$) are easily incorporated into this LP by enforcing $u_q \ge w_i r_i(x) - \alpha$ for all $i$.

However, because computing the per-probe anchor $q^{\mathrm w} = \max_{x \in \X} q(r(x))$ requires maximising a convex function over a convex set, these $2K$ global anchor subproblems may be no cheaper than front generation in the general case. In the continuous empirical demonstration (\S\ref{sec:exp}), there are $m=3$ criteria and an adaptive family of $K=7$ probes. Computing the $q^\star$ and $q^{\mathrm w}$ anchors requires $2K = 14$ global optimisation calls. The sequential lexicographic selection then evaluates up to $K=5$ active solver calls (since degenerate probes are excluded), producing the reported $19$ solver calls. We are careful not to overclaim: the core contribution of this paper remains the finite-set \emph{post hoc} selection rule, and we do not claim LexUR as a general front-free mixed-integer solver.
